% SOURCE: Chinese PDF viewer pp. 113--118. Section 6.
\section{Time-Dependent Perturbation Theory}
The Schr\"odinger equation is
\begin{equation}
 \ii\hbar\frac{\partial}{\partial t}\ket{\psi(t)}
 =\op H(t)\ket{\psi(t)}.
\end{equation}
The Hamiltonian can be written as
\begin{equation}
 \op H(t)=\op H^{(0)}+\op H^{(1)}(t),
 \qquad \op H^{(0)}\ket m=E_m\ket m.
\end{equation}
\textbf{Question.} When $t\le0$, the Hamiltonian is $\op H=\op H^{(0)}$, independent of time, and the system may be in a stationary state. Assume that at $t=0$ it is in the stationary state $\ket n$. For $t>0$, $\op H(t)$ depends on time; the system is no longer stably in the stationary state $\ket n$, but has certain probabilities of being in other stationary states. At time $t$ the state is
\begin{equation}
 \ket{\psi(t)}=\sum_m a_m(t)\ket m e^{-\ii E_mt/\hbar}.
\end{equation}
The probability of the stationary state $\ket m$ is $|a_m(t)|^2$. For a time-dependent perturbation, the quantity to be found is the probability that, after a macroscopic time ($t\to\infty$), the system has made a transition from the initial stationary state $\ket n$ to another stationary state $\ket m$:
\begin{equation}
 \lim_{t\to\infty}W_{n\to m}(t)
 =\lim_{t\to\infty}|a_m(t)|^2.
\end{equation}
\subsection*{1) Transition equations}
Substitute the expansion of $\ket{\psi(t)}$ into the Schr\"odinger equation:
\begin{equation}\begin{aligned}
 &\ii\hbar\sum_m\dot a_m\ket m e^{-\ii E_mt/\hbar}
 +\sum_m a_mE_m\ket m e^{-\ii E_mt/\hbar}\\
 &\quad=\sum_m a_m\op H^{(0)}\ket m e^{-\ii E_mt/\hbar}
 +\sum_m a_m\op H^{(1)}(t)\ket m e^{-\ii E_mt/\hbar}.
\end{aligned}\end{equation}
The second term on the left cancels the first term on the right, giving
\begin{equation}
 \ii\hbar\sum_m\dot a_m\ket m e^{-\ii E_mt/\hbar}
 =\sum_m a_m\op H^{(1)}(t)\ket m e^{-\ii E_mt/\hbar}.
\end{equation}
Multiply on the left by $\bra k$:
\begin{equation}
 \ii\hbar\dot a_ke^{-\ii E_kt/\hbar}
 =\sum_m a_m H_{km}^{(1)}(t)e^{-\ii E_mt/\hbar},
 \quad H_{km}^{(1)}(t)=\mel k{\op H^{(1)}(t)}m.
\end{equation}
Here $H_{km}^{(1)}(t)$ is the matrix element of $\op H^{(1)}(t)$ in the $\op H^{(0)}$ representation. Thus
\begin{equation}
 \boxed{\begin{aligned}
 \ii\hbar\dot a_k&=\sum_m a_m H_{km}^{(1)}(t)e^{\ii\omega_{km}t},\\
 a_k(0)&=\delta_{kn},\qquad \omega_{km}=\frac{E_k-E_m}{\hbar}.
 \end{aligned}}
\end{equation}
\subsection*{2) Perturbative solution}
Write
\begin{equation}
 a_k(t)=a_k^{(0)}(t)+a_k^{(1)}(t)+\cdots.
\end{equation}
The transition equation becomes
\begin{equation}
 \ii\hbar(\dot a_k^{(0)}+\dot a_k^{(1)}+\cdots)
 =\sum_m(a_m^{(0)}+a_m^{(1)}+\cdots)
 H_{km}^{(1)}(t)e^{\ii\omega_{km}t}.
\end{equation}
The zeroth-order equation is
\begin{equation}
 \ii\hbar\dot a_k^{(0)}(t)=0,\qquad
 a_k^{(0)}(0)=\delta_{kn},
\end{equation}
with solution
\begin{equation}
 a_k^{(0)}(t)=\delta_{kn}.
\end{equation}
The first-order equation is
\begin{equation}
 \ii\hbar\dot a_k^{(1)}(t)
 =\sum_m a_m^{(0)}(t)H_{km}^{(1)}(t)e^{\ii\omega_{km}t}
 =H_{kn}^{(1)}(t)e^{\ii\omega_{kn}t}.
\end{equation}
Its solution is
\begin{equation}
 \boxed{a_k^{(1)}(t)=\frac1{\ii\hbar}
 \int_0^t H_{kn}^{(1)}(t')e^{\ii\omega_{kn}t'}\dd t'.}
\end{equation}
For a transition from $\ket n$ to $\ket k$ ($k\ne n$), $a_k(t)\simeq a_k^{(1)}(t)$, so the transition probability is
\begin{equation}
 W_{n\to k}(t)\simeq|a_k^{(1)}(t)|^2.
\end{equation}
\subsection*{3) Periodic perturbation}
Take
\begin{equation}
 \op H^{(1)}(t)=\op A\cos\omega t
 =\op F(e^{\ii\omega t}+e^{-\ii\omega t}),
\end{equation}
where $\op A$, $\op F$ are independent of $t$. The first-order transition amplitude is
\begin{equation}\begin{aligned}
 a_m^{(1)}(t)
 &=\frac1{\ii\hbar}\int_0^t H_{mn}^{(1)}(t')e^{\ii\omega_{mn}t'}\dd t'\\
 &=\frac{F_{mn}}{\ii\hbar}\int_0^t
 \bigl[e^{\ii(\omega_{mn}+\omega)t'}
 +e^{\ii(\omega_{mn}-\omega)t'}\bigr]\dd t',
\end{aligned}\end{equation}
where $\omega_{mn}$ is the transition frequency and $\omega$ is the driving frequency. Integrating gives
\begin{equation}
 a_m^{(1)}(t)=-\frac{F_{mn}}{\hbar}\left[
 \frac{e^{\ii(\omega_{mn}+\omega)t}-1}{\omega_{mn}+\omega}
 +\frac{e^{\ii(\omega_{mn}-\omega)t}-1}{\omega_{mn}-\omega}
 \right].
\end{equation}
When $\omega\ne\pm\omega_{mn}$, $a_m^{(1)}(t)$ oscillates with $t$ and does not increase monotonically with time. When $\omega\to\pm\omega_{mn}$, the nearly resonant term dominates:
\begin{equation}
 a_m^{(1)}(t)\simeq-\frac{F_{mn}}{\hbar}
 \frac{e^{\ii(\omega_{mn}\mp\omega)t}-1}{\omega_{mn}\mp\omega}.
\end{equation}
The transition probability is
\begin{equation}
 W_{n\to m}(t)=|a_m^{(1)}(t)|^2
 \simeq\frac{4|F_{mn}|^2}{\hbar^2}
 \frac{\sin^2[(\omega_{mn}\mp\omega)t/2]}{(\omega_{mn}\mp\omega)^2}.
\end{equation}
Since
\begin{equation}
 \lim_{t\to\infty}\frac{\sin^2(xt)}{\pi t x^2}=\delta(x),
\end{equation}
it follows that
\begin{equation}
 \lim_{t\to\infty}W_{n\to m}(t)
 \simeq\frac{2\pi t|F_{mn}|^2}{\hbar^2}
 \delta(\omega_{mn}\mp\omega)
 =\frac{2\pi t|F_{mn}|^2}{\hbar}
 \delta(E_m-E_n\mp\hbar\omega).
\end{equation}
It grows linearly with $t$. The $\delta$ function expresses energy conservation, illustrated schematically below:
\begin{center}
\begin{tikzpicture}[>=Latex,scale=.9]
 \foreach \xx in {0,5.2}{\draw[thick] (\xx,0)--++(3.1,0);\draw[thick](\xx,2)--++(3.1,0);}
 \draw[->,thick] (1.55,1.9)--(1.55,.1);
 \node[right] at (1.75,1.1) {$\hbar\omega$};
 \draw[->,thick] (6.75,.1)--(6.75,1.9);
 \node[right] at (6.95,1.1) {$\hbar\omega$};
\end{tikzpicture}
\end{center}
The transition rate is
\begin{equation}
 \frac{\dd W_{n\to m}(t)}{\dd t}
 =\frac{2\pi}{\hbar^2}|F_{mn}|^2\delta(\omega_{mn}\mp\omega).
\end{equation}
A real perturbation is not a strictly monochromatic field but has a frequency range $\Delta\omega$. The total transition rate is
\begin{equation}\begin{aligned}
 \frac{\dd W_{n\to m}(t)}{\dd t}
 &=\frac{2\pi}{\hbar^2}|F_{mn}|^2
 \int_{\Delta\omega}\delta(\omega_{mn}\mp\omega)\dd\omega\\
 &=\begin{cases}
 \displaystyle\frac{2\pi}{\hbar^2}|F_{mn}|^2,
 &\Delta\omega\text{ contains }|\omega_{mn}|,\\
 0,&\Delta\omega\text{ does not contain }|\omega_{mn}|.
 \end{cases}
\end{aligned}\end{equation}
\textbf{Example.} At $t=0$ a hydrogen atom is in its ground state $\ket{100}$. A perturbation
\begin{equation}
 \op H^{(1)}(t)=-e\mathcal E_0 z\,\delta(t)
\end{equation}
is applied. Find the probability for transitions to all excited states.

\textbf{Solution.} The first-order amplitude is
\begin{equation}\begin{aligned}
 a_m^{(1)}(t)
 &=\frac1{\ii\hbar}\int_0^t H_{mn}^{(1)}(t')e^{\ii\omega_{mn}t'}\dd t'\\
 &=-\frac{e\mathcal E_0 z_{mn}}{\ii\hbar}
 \int_0^t e^{\ii\omega_{mn}t'}\delta(t')\dd t'\\
 &=-\frac{e\mathcal E_0 z_{mn}}{\ii\hbar}.
\end{aligned}\end{equation}
Here $\ket n=\ket{100}$ is the ground state and $\ket m=\ket{nlm}$ ($n\ne1$) are excited states. The probability of a transition from the ground state to any excited state is
\begin{equation}
 W=\sum_{nlm}'\left(\frac{e\mathcal E_0}{\hbar}\right)^2
 |\mel{nlm}{z}{100}|^2,
\end{equation}
where the prime means that the ground state is excluded from the sum. Since
\begin{equation}
 z=r\cos\theta,
\end{equation}
the corresponding matrix elements are
\begin{equation}
 \mel{nlm}{z}{n'l'm'}
 =\mel{nl}{r}{n'l'}\mel{lm}{\cos\theta}{l'm'}.
\end{equation}
The action of $\cos\theta$ is
\begin{equation}
 \cos\theta\ket{lm}
 =\sqrt{\frac{(l+1)^2-m^2}{(2l+1)(2l+3)}}\ket{l+1,m}
 +\sqrt{\frac{l^2-m^2}{(2l-1)(2l+1)}}\ket{l-1,m}.
\end{equation}
Therefore, the conditions for $\mel{lm}{\cos\theta}{l'm'}$ to be nonzero (the selection rules for transitions) are
\begin{equation}
 \Delta l=l-l'=\pm1.
\end{equation}
Thus $\mel{100}{z}{100}=0$, and
\begin{equation}\begin{aligned}
 W&=\sum_{nlm}\left(\frac{e\mathcal E_0}{\hbar}\right)^2
 |\mel{nlm}{z}{100}|^2\\
 &=\left(\frac{e\mathcal E_0}{\hbar}\right)^2
 \sum_{nlm}\mel{100}{z}{nlm}\mel{nlm}{z}{100}\\
 &=\left(\frac{e\mathcal E_0}{\hbar}\right)^2\mel{100}{z^2}{100}.
\end{aligned}\end{equation}
Since
\begin{equation}
 \mel{100}{z^2}{100}=a_0^2,
\end{equation}
where $a_0$ is the Bohr radius, the result is
\begin{equation}
 \boxed{W=\left(\frac{e\mathcal E_0 a_0}{\hbar}\right)^2}.
\end{equation}
The probability for the hydrogen atom to remain in the ground state is therefore
\begin{equation}
 \boxed{1-W=1-\left(\frac{e\mathcal E_0 a_0}{\hbar}\right)^2}.
\end{equation}

\interactiveguide{IC-04-03}{Driven two-level resonance}{After the periodic-drive transition calculation, scan the drive frequency and duration. Examine resonant versus off-resonant transition probabilities and the finite-time linewidth; compare weak-driving predictions with the exact two-level model.}
